\chapter{Git Integration}
\label{ch:git}

\section{Git Metadata Commands}

The template provides four macros for embedding repository metadata in the
document, all defined in \texttt{omnilatex-utils.sty}:

\begin{description}
  \item[\textbackslash GitRefName] The current branch or tag name.
        \begin{minted}{latex}
        Compiled from branch: \GitRefName
        \end{minted}

  \item[\textbackslash GitShortSHA] The abbreviated commit hash (8 characters).
        \begin{minted}{latex}
        Commit: \GitShortSHA
        \end{minted}

  \item[\textbackslash GitLongSHA] The full 40-character commit hash.
        \begin{minted}{latex}
        Full SHA: \GitLongSHA
        \end{minted}

  \item[\textbackslash BuildDate] The date and time of compilation.
        \begin{minted}{latex}
        Built on \BuildDate.
        \end{minted}
\end{description}

All Git macros require \texttt{shell-escape} to be enabled, which is the
default in all build modes.

\section{Verification Box}

\cs{GitVerificationBox} renders a formatted box suitable for colophon pages:

\begin{minted}{latex}
\GitVerificationBox
\end{minted}

The box displays:
\begin{itemize}
  \item The full and short commit SHA.
  \item The branch or tag name.
  \item A verification command the reader can run to confirm the build.
\end{itemize}

\section{Build Date}

\cs{BuildDate} integrates with \texttt{datetime2} (\cs{DTMnow}) to produce
a locale-aware timestamp:

\begin{minted}{latex}
This document was compiled on \BuildDate.
\end{minted}

The format respects the document's language and date settings.

\section{Practical Use}

A common pattern is to place Git metadata on the colophon or back-matter
page:

\begin{minted}{latex}
\chapter*{Colophon}
This document was built from commit \GitShortSHA{} on branch
\GitRefName{} at \BuildDate.  To verify the build:
\begin{minted}[fontsize=\small]{bash}
git clone <repo-url>
git checkout \GitLongSHA
BUILD_MODE=prod python build.py
\end{minted}
\GitVerificationBox
\end{minted}

This ensures full reproducibility: any reader can check out the exact
revision and reproduce the PDF.
